% Generated draft: requires mathematical and source review.
\documentclass[11pt]{article}
\usepackage[margin=25mm]{geometry}
\usepackage{fontspec}
\setmainfont{Latin Modern Roman}
\usepackage{amsmath,amssymb,mathtools,unicode-math}
\setmathfont{Latin Modern Math}
\usepackage{enumitem,needspace,xurl,hyperref}
\setlength{\parindent}{0pt}
\setlength{\parskip}{5pt}
\setlength{\emergencystretch}{3em}
\newcommand{\N}{\mathbb{N}}
\newcommand{\Z}{\mathbb{Z}}
\newcommand{\Q}{\mathbb{Q}}
\newcommand{\R}{\mathbb{R}}
\newcommand{\C}{\mathbb{C}}
\newcommand{\F}{\mathbb{F}}
\newcommand{\A}{\mathbb{A}}
\newcommand{\PP}{\mathbb{P}}
\newcommand{\E}{\mathbb{E}}
\newcommand{\eps}{\varepsilon}
\newcommand{\dd}{\,\mathrm{d}}
\begin{document}
\section*{UnsolvedMath ID: 700689}
\textbf{Source problem code:} catalogue entry
\subsection{Definitions and mathematical statement}
Let \(G\) be a group, and let \(a\in G\) have order \(p\), where \(p\) is a prime. For \(g\in G\), write \(a^g=g^{-1}ag\), and define
\[
C_G(a)=\{x\in G:xa=ax\},\qquad H_g=\langle a,a^g\rangle,\qquad \mathcal H=\{H_g:g\in G\}.
\]
A group \(H\) is soluble if its derived series, defined by \(H^{(0)}=H\) and \(H^{(n+1)}=[H^{(n)},H^{(n)}]\), reaches the trivial subgroup after finitely many steps. A group is locally finite if every subgroup generated by finitely many elements is finite.

Interpret ``almost all'' in the supplied statement as ``all but finitely many distinct subgroups in \(\mathcal H\).'' The source record does not explicitly define this phrase, so this interpretation requires confirmation against the source.

Under this interpretation, the question is: if
\[
|C_G(a)|<\infty,\qquad (\forall g\in G)\ |H_g|<\infty,\qquad
\bigl|\{H\in\mathcal H:H\text{ is not soluble}\}\bigr|<\infty,
\]
must \(G\) be locally finite? Equivalently, must
\[
(\forall F\subseteq G)\bigl(|F|<\infty\Longrightarrow |\langle F\rangle|<\infty\bigr)
\]
hold?

The supplied source describes this problem as closely connected with Kourovka Notebook Problem 6.56 and reports positive answers for several important partial cases obtained by the proposer. Neither the statement of Problem 6.56 nor the hypotheses of those partial cases are supplied.
\subsection{Short English statement}
Suppose a group contains an element of prime order with finite centralizer. Every subgroup generated by that element and one of its conjugates is finite, and almost all such subgroups are soluble. Must every finitely generated subgroup of the group be finite? Here ``almost all'' is explicitly interpreted as allowing only finitely many exceptional distinct subgroups; the supplied record leaves that phrase undefined.
\subsection{Sources}
\begin{itemize}
\item Supplied local UnsolvedMath registry record, permanent ID \(700689\), exact source problem code \texttt{catalogue entry}, titled \emph{Kourovka Notebook Problem 10.74}. The registry labels the problem \emph{open}; that label is supplied metadata, not an independently verified mathematical assessment.
\item \emph{The Kourovka Notebook (Unsolved Problems in Group Theory)}, original Problem 10.74, 10th issue (1986), proposed by V. P. Shunkov. The supplied attribution identifies the 21st issue edition (2026), \texttt{arXiv:1401.0300v48}; the collection metadata names E. I. Khukhro and V. D. Mazurov as editors. These bibliographic details have not been independently verified.
\item V. P. Shunkov, contribution cited in \emph{Abstracts on Group Theory of the Mal'cev Int. Conf. on Algebra}, Novosibirsk, 1989, p. 145 (Russian). The supplied statement attributes positive solutions of several important partial cases to the proposer and cites this reference, without specifying those cases or a contribution title.
\item The supplied statement cross-references Kourovka Notebook Problem 6.56 as closely connected; its formulation is not included in the record.
\end{itemize}
\end{document}
